\section{Conclusion}
\label{sec:conclusion}

This study conducted a source-stratified empirical benchmark evaluating closed-set fruit-quality classification and post-hoc out-of-distribution detection across $N = 13,465$ photographic records under domain, semantic species, and dataset-reported treatment shifts. Our findings provide empirical answers to the three core research questions:
\begin{enumerate}
\setlength{\itemsep}{0pt}\setlength{\parsep}{0pt}\setlength{\parskip}{0pt}\setlength{\topsep}{1pt}\setlength{\partopsep}{0pt}
\item \textbf{Resolution of RQ1 (External Domain Robustness):} Evaluating a frozen MobileNetV2 classifier on external repository imagery reveals a marked 28.56 percentage-point degradation in top-1 accuracy (from 96.92\% [94.71\%, 98.68\%] on laboratory test data to 68.36\% [67.20\%, 69.40\%] on external knowns). This degradation is accompanied by a pronounced class-wise failure on Fresh Apple (32.12\% accuracy; 310/965 correct), where natural surface striations, lenticel spots, and specularity are observed to coincide with decay misclassifications (393 Rotten Apple, 255 Rotten Orange), which we present strictly as an untested observational hypothesis.
\item \textbf{Resolution of RQ2 (Comparative Out-of-Distribution Discrimination):} Under source-stratified semantic species shifts, Energy scoring (AUROC = 81.57\% on FruitVision, 67.52\% on Sultana) and Ledoit--Wolf Mahalanobis distance (81.41\% and 68.15\%) achieve higher AUROC point estimates than baseline Maximum Softmax Probability (72.42\% and 59.77\%). However, on dataset-reported chemical preservative shift (Cohort B), all evaluated post-hoc scoring algorithms exhibit weak or near-chance discrimination (Energy AUROC = 53.13\%, Mahalanobis = 47.91\%, FPR@95 $\ge 91.13\%$). Post-hoc scoring methods on standard RGB inputs do not reliably separate treatment-labelled images from natural produce.
\item \textbf{Resolution of RQ3 (Operational Feasibility and Rejection Trade-offs):} Prospective validation-calibrated gating ($\tau_{95}$) achieves partial screening on pure semantic species shift in Cohort A (Dual-Gate UFAR = 16.07\%) and on combined species-and-treatment shift in Cohort C (Dual-Gate UFAR = 9.54\%), but incurs a substantial in-distribution abstention penalty, rejecting 42.43\% [41.10\%, 43.76\%] of external known-class image records (2,270/5,350) and retaining automated decisions on only 57.57\% of image records. Crucially, 42.83\% [41.26\%, 44.34\%] of dataset-reported treatment-labelled image records (821/1,917) are accepted into the classifier's Fresh-class output labels.
\end{enumerate}

\noindent This investigation demonstrates the critical necessity of source-stratified evaluation in agricultural AI and exposes the visual failure boundaries of standalone RGB computer vision. Future work could evaluate whether multi-modal architectures integrating optical imaging with internal spectroscopy or chemical headspace sensing improve screening under these condition shifts.
